\subsection{Qudit Clifford Group}

The Qudit Clifford group $\mathcal{C}_q$ is the normalizer of the $n$-qudit Pauli group $\mathcal{P}_q^{\otimes n}$ within the unitary group $U(q^n)$. It is defined as the set of operations that map Pauli operators to Pauli operators under conjugation \cite{Gottesman1999}:
\begin{equation}
    \mathcal{C}_q = \{ U \in U(q^n) \mid U \mathcal{P}_q^{\otimes n} U^\dagger = \mathcal{P}_q^{\otimes n} \}
\end{equation}
This property is central to the Gottesman-Knill theorem, which states that quantum circuits composed solely of Clifford gates and Pauli measurements can be efficiently simulated on a classical computer \cite{Gottesman1998}.

For a prime dimension $q$, the Clifford group is generated by three elementary gates: the Hadamard gate $H$, the Phase gate $S$, and the two-qudit $SUM$ gate (generalizing the qubit CNOT) \cite{Hostens2005}.

\subsubsection{Generators and Conjugation Rules}
The action of these generators on the computational basis and their conjugation effects on the Pauli operators are defined as follows:

\textbf{1. The Hadamard Gate ($H$)} \\
The Hadamard gate implements a Discrete Fourier Transform over $\mathbb{Z}_q$, changing the basis from computational ($Z$) to conjugate ($X$):
\begin{equation}
    H |j\rangle = \frac{1}{\sqrt{q}} \sum_{k=0}^{q-1} \omega^{jk} |k\rangle
\end{equation}
For arbitrary $q$, its conjugation action rotates the Pauli operators, mapping $X$ to $Z$ and $Z$ to $X^{q-1}$:
\begin{align}
    H X H^\dagger &= Z, \\
    H Z H^\dagger &= X^\dagger = X^{q-1} = X^{-1}
\end{align}

\textbf{2. The Phase Gate ($S$)} \\
The Phase gate applies a quadratic phase to the basis states. For odd prime $q$ (such as qutrits, $q=3$), it is defined as:
\begin{equation}
    S |j\rangle = \omega^{j(j-1)/2} |j\rangle
\end{equation}
For $q=3$ this gives $S=\mathrm{diag}(1,1,\omega)$. The \texttt{sdim} simulator calls this gate \texttt{P}.
This gate preserves $Z$ operators but mixes $X$ and $Z$:
\begin{align}
    S X S^\dagger &= X Z, \\
    S Z S^\dagger &= Z
\end{align}

\textbf{3. The Sum Gate ($SUM$)} \\
The $SUM$ gate is the canonical two-qudit entangling gate, generalizing the CNOT. It adds the state of the control qudit ($c$) to the target qudit ($t$) modulo $q$:
\begin{equation}
    SUM (|j\rangle_c \otimes |k\rangle_t)= |j\rangle_c \otimes |(j+k) \bmod q\rangle_t.
\end{equation}
The conjugation rules describe the propagation of errors between qudits:
\begin{align}
    SUM (X \otimes I) SUM^\dagger &= X \otimes X, \\
    SUM (I \otimes X) SUM^\dagger &= I \otimes X, \\
    SUM (Z \otimes I) SUM^\dagger &= Z \otimes I, \\
    SUM (I \otimes Z) SUM^\dagger &= Z^{-1} \otimes Z
\end{align}
Note the appearance of $Z^{-1}$ on the control qudit when a $Z$ error propagates from the target, a distinct feature of qudit stabilizer logic ($q>2$) compared to qubits.

\textbf{4. The Controlled-Phase Gate ($CZ$)} \\
The controlled-phase gate used in the logical $\bar S$ gate follows from the generators. It is diagonal:
\begin{equation}
    CZ = (I\otimes H)\,SUM\,(I\otimes H^\dagger), \qquad
    CZ (|j\rangle \otimes |k\rangle) = \omega^{jk}\, |j\rangle \otimes |k\rangle,
\end{equation}
so it leaves $Z\otimes I$ and $I\otimes Z$ unchanged and maps $X\otimes I \mapsto X\otimes Z$ and $I\otimes X \mapsto Z\otimes X$.
